% =====================================================================
% Artículo IEEE (conferencia) - Sistema de traducción automática
% inglés -> latín/griego con ajuste de hiperparámetros mediante
% MAX-MIN Ant System.
%
% Compilar: pdflatex main -> bibtex main -> pdflatex main -> pdflatex main
% (o subir la carpeta paper/ a Overleaf).
%
% Todo dato aún no disponible está marcado con \pendiente{...}, que se
% imprime en rojo. Antes de entregar: buscar "\pendiente" y eliminarlos.
% =====================================================================
\documentclass[conference]{IEEEtran}

\usepackage[T1]{fontenc}
\usepackage[utf8]{inputenc}
\usepackage[spanish,es-noshorthands,es-tabla]{babel}
\usepackage{cite}
\usepackage{amsmath,amssymb}
\usepackage{graphicx}
\usepackage{booktabs}
\usepackage{multirow}
\usepackage{algorithm}
\usepackage{algpseudocode}
\usepackage{xcolor}
\usepackage{url}

% --- Utilidades -------------------------------------------------------
\newcommand{\pendiente}[1]{\textcolor{red}{[PENDIENTE: #1]}}
\newcommand{\chrf}{chrF}
\newcommand{\mmas}{$\mathcal{MAX}$--$\mathcal{MIN}$ Ant System}

% Incluye la figura si existe; si no, dibuja un recuadro de marcador.
\newcommand{\figuraopcional}[2]{%
  \IfFileExists{#1}{\includegraphics[width=#2]{#1}}{%
    \fbox{\parbox[c][3.5cm][c]{0.9#2}{\centering\pendiente{figura \texttt{\detokenize{#1}}}}}}}

\floatname{algorithm}{Algoritmo}
\renewcommand{\algorithmicrequire}{\textbf{Entrada:}}
\renewcommand{\algorithmicensure}{\textbf{Salida:}}

\begin{document}

\title{Traducción automática del inglés al latín y al griego antiguo con
ajuste de hiperparámetros mediante colonia de hormigas \mmas}

% TODO: completar nombres, afiliación y correos de los 4 integrantes.
\author{
\IEEEauthorblockN{Integrante 1, Integrante 2, Integrante 3, Integrante 4}
\IEEEauthorblockA{\pendiente{Departamento, Universidad}\\
\pendiente{Ciudad, País}\\
\pendiente{correos}}
}

\maketitle

% =====================================================================
\begin{abstract}
La traducción automática hacia lenguas clásicas es un problema de bajos
recursos con morfología flexiva rica, en el que la elección de
hiperparámetros de un modelo neuronal influye de forma importante en la
calidad. Este trabajo presenta un sistema de traducción del inglés al
latín y al griego antiguo basado en un Transformer codificador--decodificador
entrenado desde cero, cuyos hiperparámetros se ajustan con \mmas{} (MMAS),
un algoritmo de optimización por colonia de hormigas. El espacio de
búsqueda se discretiza en un grafo con nueve hiperparámetros, y la aptitud
de cada hormiga es el \chrf{} en validación tras un entrenamiento de
presupuesto fijo. MMAS se compara con una configuración por defecto y con
una búsqueda aleatoria de igual presupuesto. En el conjunto de prueba de
latín, la configuración encontrada por MMAS obtiene \pendiente{X} de
\chrf{} frente a \pendiente{Y} de la búsqueda aleatoria y \pendiente{Z}
de la configuración por defecto; transferida al griego, obtiene
\pendiente{W}. \pendiente{Conclusión principal en una frase.}
\end{abstract}

\begin{IEEEkeywords}
traducción automática neuronal, optimización por colonia de hormigas,
MAX--MIN Ant System, ajuste de hiperparámetros, latín, griego antiguo,
Transformer
\end{IEEEkeywords}

% =====================================================================
\section{Introducción}

La traducción automática neuronal basada en la arquitectura Transformer
\cite{vaswani2017attention} es hoy el enfoque dominante para pares de
lenguas con abundantes datos paralelos. Las lenguas clásicas, como el
latín y el griego antiguo, presentan en cambio dos dificultades: los
corpus paralelos son pequeños o heterogéneos, y su morfología flexiva hace
que una misma palabra aparezca en muchas formas distintas. En este
contexto, la calidad del modelo depende en gran medida de decisiones de
diseño como la profundidad, el tamaño de las representaciones, la
regularización y la tasa de aprendizaje.

Ajustar estos hiperparámetros es costoso: cada configuración candidata
exige entrenar un modelo. La búsqueda manual depende de la experiencia
del investigador, y la búsqueda en rejilla crece exponencialmente con el
número de hiperparámetros. La búsqueda aleatoria es una alternativa
sólida y difícil de superar \cite{bergstra2012random}, pero no aprovecha
la información de las evaluaciones anteriores.

Los algoritmos de optimización por colonia de hormigas (ACO)
\cite{dorigo2004aco} construyen soluciones de forma incremental guiados
por un rastro de feromona que acumula la experiencia de la colonia. En
particular, \mmas{} (MMAS) \cite{stutzle2000mmas} limita la feromona a un
intervalo acotado para evitar el estancamiento prematuro de la búsqueda,
lo que lo hace atractivo cuando el número de evaluaciones es pequeño.

Las contribuciones de este trabajo son:
\begin{itemize}
  \item Un sistema de traducción del inglés al latín y al griego antiguo
  basado en un Transformer entrenado desde cero con tokenizadores de
  subpalabras propios.
  \item La formulación del ajuste de hiperparámetros como la construcción
  de un camino en un grafo discretizado, resuelto con MMAS.
  \item Una comparación con igual presupuesto entre MMAS, la búsqueda
  aleatoria y una configuración por defecto, con pruebas de significancia
  estadística.
  \item Un análisis de la transferencia de la configuración ajustada en
  latín al griego antiguo.
\end{itemize}

% =====================================================================
\section{Trabajo relacionado}

\subsection{Traducción automática neuronal}
El Transformer \cite{vaswani2017attention} sustituyó a las redes
recurrentes mediante mecanismos de atención, y es la base de herramientas
como Marian \cite{junczys2018marian}, cuya arquitectura está disponible en
la biblioteca Transformers \cite{wolf2020transformers}. La segmentación en
subpalabras mediante \emph{byte pair encoding} (BPE)
\cite{sennrich2016bpe}, implementada en SentencePiece
\cite{kudo2018sentencepiece}, permite modelar vocabularios abiertos, lo
que resulta especialmente útil en lenguas de morfología rica.

\subsection{Optimización por colonia de hormigas}
ACO \cite{dorigo2004aco} se inspira en el comportamiento de forrajeo de
las hormigas, que depositan feromona en los caminos y prefieren los
caminos con más feromona. MMAS \cite{stutzle2000mmas} introduce dos
modificaciones al Ant System original: solo la mejor hormiga deposita
feromona, y la feromona se acota en $[\tau_{\min}, \tau_{\max}]$; con
ellas se encuentra entre los algoritmos de mejor desempeño en problemas
como el del agente viajero y el de asignación cuadrática. Para dominios
continuos, ACO$_\mathbb{R}$ \cite{socha2008acor} representa la feromona
como un archivo de soluciones y muestrea de mezclas de gaussianas, y
ACO$_\text{MV}$ \cite{liao2014acomv} lo extiende a problemas con
variables continuas, ordinales y categóricas.

\subsection{ACO para el diseño de redes neuronales}
DeepSwarm \cite{byla2019deepswarm} aplica ACO a la búsqueda de
arquitecturas de redes convolucionales: cada hormiga recorre un grafo de
capas candidatas, el camino completo define una red que se evalúa, y la
mejor hormiga refuerza su camino. Nuestro trabajo adopta la misma
representación de grafo de opciones, pero la aplica al ajuste de
hiperparámetros de un modelo de traducción.

\subsection{Evaluación de la traducción}
BLEU \cite{papineni2002bleu} mide la coincidencia de $n$-gramas de
palabras, mientras que \chrf{} \cite{popovic2015chrf} usa $n$-gramas de
caracteres y otorga crédito parcial a formas flexionadas distintas de la
referencia. Seguimos las recomendaciones de sacreBLEU
\cite{post2018sacrebleu} para reportar puntuaciones reproducibles, y el
remuestreo \emph{bootstrap} pareado \cite{koehn2004bootstrap} para la
significancia estadística.

% =====================================================================
\section{Datos y preprocesamiento}

\subsection{Corpus}
Se usan dos corpus paralelos públicos: latín--inglés
\cite{dataset_latin} y griego antiguo--inglés \cite{dataset_griego}. La
Tabla~\ref{tab:corpus} resume su tamaño tras el preprocesamiento.

\subsection{Normalización y partición}
Cada oración se normaliza a la forma Unicode NFC, se colapsan los espacios
en blanco y se convierte a minúsculas. Los pares nulos y duplicados se
eliminan \emph{antes} de la partición, para evitar fuga de información
entre particiones. Cada corpus se divide aleatoriamente (semilla 42) en
80\,\% entrenamiento, 10\,\% validación y 10\,\% prueba.

\subsection{Tokenización}
Para cada par de lenguas se entrena un tokenizador SentencePiece BPE
conjunto, sobre ambos lados del corpus de entrenamiento, con cobertura de
caracteres completa, sin normalización adicional y con \emph{byte
fallback} para evitar símbolos desconocidos. El tamaño del vocabulario se
eligió comparando la fertilidad (tokens por palabra) en validación para
varios tamaños (Tabla~\ref{tab:tokenizadores}): 8\,000 para latín y
16\,000 para griego (este último entrenado sobre una muestra de 300\,000
oraciones). Los símbolos \texttt{<2en>}, \texttt{<2la>} y \texttt{<2el>}
se reservan como etiquetas de dirección indivisibles para una futura
versión bidireccional.

\subsection{Filtrado por longitud}
Se descartan los pares con más de 160 tokens (latín) o 128 tokens
(griego) en cualquiera de los lados. En entrenamiento se descartan además
los pares cuya razón de longitudes supera 3, como indicio de mala
alineación; este filtro no se aplica a validación ni a prueba, para no
facilitar artificialmente la evaluación.

\begin{table}[t]
\centering
\caption{Tamaño de los corpus (pares de oraciones) tras deduplicación,
partición y filtrado por longitud}
\label{tab:corpus}
\begin{tabular}{lrrr}
\toprule
Corpus & Entrenamiento & Validación & Prueba \\
\midrule
Latín--inglés  & 78\,463 & 9\,821 & 9\,812 \\
Griego--inglés & 1\,101\,526 & 139\,093 & 138\,900 \\
\midrule
Griego (submuestra) & 300\,000 & -- & -- \\
\bottomrule
\end{tabular}
\end{table}

% NOTA: cifras preliminares, obtenidas con tokenizadores sin etiquetas
% reservadas. Recalcular tras reentrenar (las diferencias deben ser mínimas).
\begin{table}[t]
\centering
\caption{Fertilidad (tokens por palabra) en validación según el tamaño
del vocabulario. En negrita, el tamaño elegido. \pendiente{recalcular}}
\label{tab:tokenizadores}
\begin{tabular}{llrr}
\toprule
Par & Vocabulario & Lengua clásica & Inglés \\
\midrule
\multirow{3}{*}{la--en} & 4\,000  & 2.012 & 1.526 \\
                        & \textbf{8\,000}  & \textbf{1.756} & \textbf{1.365} \\
                        & 16\,000 & 1.551 & 1.260 \\
\midrule
\multirow{3}{*}{el--en} & 8\,000  & 1.565 & 1.364 \\
                        & \textbf{16\,000} & \textbf{1.346} & \textbf{1.233} \\
                        & 32\,000 & 1.220 & 1.166 \\
\bottomrule
\end{tabular}
\end{table}

% =====================================================================
\section{Modelo de traducción}

Se usa la arquitectura Marian \cite{junczys2018marian}, un Transformer
codificador--decodificador \cite{vaswani2017attention}, en su
implementación de la biblioteca Transformers \cite{wolf2020transformers},
\emph{inicializada aleatoriamente} (sin preentrenamiento). Las matrices de
\emph{embeddings} del codificador, del decodificador y de la capa de
salida se comparten, dado que el vocabulario es conjunto. Los
identificadores especiales coinciden con los del tokenizador
(desconocido 0, inicio 1, fin 2, relleno 3).

El entrenamiento usa un bucle propio con AdamW, calendario de raíz
cuadrada inversa con calentamiento, precisión mixta (fp16) y lotes
formados por número de tokens (aproximadamente 8\,000 tokens por lote).
Durante la búsqueda, cada candidato se entrena con un \emph{presupuesto
fijo de tokens objetivo}, equivalente a unas dos épocas del corpus de
latín, con la misma semilla para todos los candidatos. Un presupuesto en
tokens, a diferencia de uno en tiempo o en pasos, no favorece a los
modelos más rápidos ni a los de lotes más grandes.

% =====================================================================
\section{Ajuste de hiperparámetros con MMAS}

\subsection{Espacio de búsqueda como grafo}
El espacio de búsqueda se discretiza en nueve hiperparámetros
(Tabla~\ref{tab:espacio}). Se representa como un grafo por capas: la capa
$h$ contiene un nodo por cada valor $v \in V_h$ del hiperparámetro $h$.
Una hormiga construye una configuración recorriendo las capas en orden y
eligiendo un nodo en cada una. El espacio contiene
$3 \cdot 2 \cdot 5 \cdot 5 \cdot 2 \cdot 4 \cdot 5 \cdot 4 \cdot 3 = 144\,000$
configuraciones, todas válidas, pues cada dimensión del modelo es
divisible entre cada número de cabezas.

\begin{table}[t]
\centering
\caption{Espacio de búsqueda discretizado}
\label{tab:espacio}
\begin{tabular}{ll}
\toprule
Hiperparámetro & Valores \\
\midrule
Dimensión del modelo $d_\text{model}$ & 256, 384, 512 \\
Cabezas de atención & 4, 8 \\
Capas del codificador & 2, 3, 4, 5, 6 \\
Capas del decodificador & 2, 3, 4, 5, 6 \\
Multiplicador de la red \emph{feed-forward} & 2, 4 \\
\emph{Dropout} & 0.0, 0.1, 0.2, 0.3 \\
Tasa de aprendizaje & $10^{-4}$, $2{\cdot}10^{-4}$, $3{\cdot}10^{-4}$, $5{\cdot}10^{-4}$, $10^{-3}$ \\
Pasos de calentamiento & 500, 1\,000, 2\,000, 4\,000 \\
Suavizado de etiquetas & 0.0, 0.1, 0.2 \\
\bottomrule
\end{tabular}
\end{table}

\subsection{Función de aptitud}
La aptitud $f(s) \in [0, 1]$ de una configuración $s$ es el \chrf{}
(dividido entre 100) obtenido con decodificación voraz sobre una muestra
fija de 1\,000 pares de validación, tras entrenar con el presupuesto fijo
de tokens. Las configuraciones con más de 40 millones de parámetros, o
que agotan la memoria de la GPU, reciben aptitud 0. Todas las
evaluaciones se almacenan en una caché, de modo que un camino repetido no
se vuelve a entrenar.

\subsection{Construcción de soluciones}
Sea $\tau_{h,v}$ la feromona del valor $v$ en la capa $h$. La hormiga
elige el valor $v$ con probabilidad
\begin{equation}
  p_{h,v} = \frac{\tau_{h,v}^{\alpha}}{\sum_{u \in V_h} \tau_{h,u}^{\alpha}},
  \label{eq:transicion}
\end{equation}
con $\alpha = 1$. No se usa información heurística ($\beta = 0$), para no
introducir conocimiento previo que favorezca a MMAS frente a la búsqueda
aleatoria. La feromona inicial es uniforme, por lo que la primera
iteración de MMAS equivale a un muestreo aleatorio.

\subsection{Actualización de la feromona}
Tras evaluar las $m$ hormigas de la iteración, la feromona se evapora y
solo la mejor hormiga $s^{\text{best}}$ deposita:
\begin{equation}
  \tau_{h,v} \leftarrow (1 - \rho)\,\tau_{h,v} +
  \begin{cases}
    f(s^{\text{best}}) & \text{si } v \in s^{\text{best}},\\
    0 & \text{en otro caso,}
  \end{cases}
  \label{eq:actualizacion}
\end{equation}
donde $s^{\text{best}}$ es la mejor hormiga de la iteración, y en las
últimas iteraciones la mejor encontrada hasta el momento
$s^{\text{gb}}$. Finalmente, la feromona se acota:
\begin{equation}
  \tau_{h,v} \leftarrow \min\bigl(\tau_{\max}, \max(\tau_{\min}, \tau_{h,v})\bigr),
  \quad
  \tau_{\max} = \frac{f(s^{\text{gb}})}{\rho},
  \quad
  \tau_{\min} = \frac{\tau_{\max}}{a}.
  \label{eq:cotas}
\end{equation}
Las cotas impiden que la probabilidad de un valor llegue a cero, lo que
mantiene la exploración con un número reducido de evaluaciones. El
Algoritmo~\ref{alg:mmas} resume el procedimiento.

\begin{algorithm}[t]
\caption{MMAS para el ajuste de hiperparámetros}
\label{alg:mmas}
\begin{algorithmic}[1]
\Require capas $\{V_h\}$, hormigas $m$, iteraciones $T$, $\alpha$, $\rho$, $a$
\Ensure mejor configuración $s^{\text{gb}}$
\State $\tau_{h,v} \gets 1$ para todo $h, v$; \quad $s^{\text{gb}} \gets \varnothing$
\For{$t = 1$ \textbf{to} $T$}
  \For{$k = 1$ \textbf{to} $m$}
    \State construir $s_k$ eligiendo $v \in V_h$ con la Ec.~\eqref{eq:transicion}
    \If{$s_k$ está en la caché}
       \State reutilizar $f(s_k)$ y construir una hormiga adicional
    \Else
       \State entrenar y evaluar: $f(s_k)$; guardar en la caché
    \EndIf
  \EndFor
  \State actualizar $s^{\text{gb}}$; elegir $s^{\text{best}}$ (iteración o global)
  \State evaporar y depositar con la Ec.~\eqref{eq:actualizacion}
  \State acotar con la Ec.~\eqref{eq:cotas}
\EndFor
\State \Return $s^{\text{gb}}$
\end{algorithmic}
\end{algorithm}

% =====================================================================
\section{Configuración experimental}

\subsection{Métodos comparados}
\begin{itemize}
  \item \textbf{Por defecto}: $d_\text{model}=512$, 8 cabezas, 3+3 capas,
  multiplicador 4, \emph{dropout} 0.1, tasa de aprendizaje
  $5{\cdot}10^{-4}$, 2\,000 pasos de calentamiento y suavizado 0.1.
  \item \textbf{Búsqueda aleatoria}: muestreo uniforme sobre la misma
  rejilla de la Tabla~\ref{tab:espacio}.
  \item \textbf{MMAS}: $m = 8$ hormigas, $T = 5$ iteraciones,
  $\alpha = 1$, $\beta = 0$, $\rho = 0.2$, $a = $ \pendiente{valor de $a$}.
\end{itemize}
La búsqueda aleatoria y MMAS disponen del mismo presupuesto de 40
evaluaciones, la misma función de aptitud, las mismas semillas, el mismo
tope de parámetros y la misma muestra de validación.

\subsection{Búsqueda sobre latín y transferencia al griego}
La búsqueda se ejecuta sobre el par inglés$\to$latín, cuyo tamaño permite
decenas de evaluaciones. La mejor configuración se transfiere sin cambios
al par inglés$\to$griego, entrenado sobre una submuestra fija de 300\,000
pares por un máximo de 6 épocas.

\subsection{Entrenamiento final y evaluación}
La configuración por defecto, la mejor de la búsqueda aleatoria y la
mejor de MMAS se entrenan hasta convergencia, con parada temprana según
el \chrf{} sobre la validación completa. El conjunto de prueba se usa una
sola vez, al final. La decodificación usa búsqueda en haz de tamaño 5 y
penalización de longitud 1.0. Se reportan \chrf{} (métrica principal) y
BLEU con sacreBLEU \cite{post2018sacrebleu} sobre el texto destokenizado,
y la significancia de las diferencias se estima con \emph{bootstrap}
pareado \cite{koehn2004bootstrap}.

\subsection{Infraestructura}
MMAS y el entrenamiento final de latín se ejecutan en una GPU NVIDIA RTX
4060 Laptop (8\,GB); la búsqueda aleatoria y los modelos griegos, en GPU
NVIDIA T4 de Google Colab. Cada resultado registra el \emph{commit} del
código, la huella del tokenizador, la GPU y las versiones de las
bibliotecas. \pendiente{versiones de PyTorch, Transformers, SentencePiece
y sacreBLEU}.

% =====================================================================
\section{Resultados}

\subsection{Búsqueda de hiperparámetros}
La Fig.~\ref{fig:convergencia} muestra la mejor aptitud alcanzada en
función del número de evaluaciones para MMAS y la búsqueda aleatoria, y
la Fig.~\ref{fig:feromona} la evolución de la feromona.
\pendiente{describir la convergencia y la configuración encontrada}.

\begin{figure}[t]
\centering
\figuraopcional{figuras/convergencia.pdf}{\columnwidth}
\caption{Mejor aptitud (\chrf{} en validación) frente al número de
evaluaciones en inglés$\to$latín.}
\label{fig:convergencia}
\end{figure}

\begin{figure}[t]
\centering
\figuraopcional{figuras/feromona.pdf}{\columnwidth}
\caption{Distribución de la feromona por hiperparámetro al final de cada
iteración de MMAS.}
\label{fig:feromona}
\end{figure}

\begin{table}[t]
\centering
\caption{Mejores configuraciones encontradas}
\label{tab:configuraciones}
\begin{tabular}{lccc}
\toprule
Hiperparámetro & Por defecto & Aleatoria & MMAS \\
\midrule
$d_\text{model}$ & 512 & -- & -- \\
Cabezas & 8 & -- & -- \\
Capas cod./dec. & 3/3 & -- & -- \\
Multiplicador FFN & 4 & -- & -- \\
\emph{Dropout} & 0.1 & -- & -- \\
Tasa de aprendizaje & $5{\cdot}10^{-4}$ & -- & -- \\
Calentamiento & 2\,000 & -- & -- \\
Suavizado & 0.1 & -- & -- \\
Parámetros (M) & -- & -- & -- \\
\chrf{} validación (búsqueda) & -- & -- & -- \\
\bottomrule
\end{tabular}
\end{table}

\subsection{Evaluación en prueba}
\pendiente{describir la Tabla~\ref{tab:resultados}}.

\begin{table}[t]
\centering
\caption{Resultados en el conjunto de prueba (haz 5). $^\dagger$:
diferencia significativa respecto a MMAS ($p < 0.05$, \emph{bootstrap}
pareado)}
\label{tab:resultados}
\begin{tabular}{llrr}
\toprule
Par & Configuración & \chrf{} & BLEU \\
\midrule
\multirow{3}{*}{en$\to$la} & Por defecto & -- & -- \\
                           & Aleatoria   & -- & -- \\
                           & MMAS        & -- & -- \\
\midrule
\multirow{2}{*}{en$\to$el} & Por defecto & -- & -- \\
                           & MMAS (transferida) & -- & -- \\
\bottomrule
\end{tabular}
\end{table}

\subsection{Análisis cualitativo}
\pendiente{unas 10 traducciones de ejemplo y análisis de errores:
desinencias, orden de palabras, nombres propios}.

% =====================================================================
\section{Discusión}

\pendiente{interpretar los resultados}.

\textbf{Limitaciones.} El presupuesto corto de la función de aptitud
favorece configuraciones que aprenden rápido, que no necesariamente son
las mejores a convergencia; por ello los finalistas se entrenan hasta
convergencia antes de compararlos. La rejilla se eligió a mano, y con
40 evaluaciones solo se explora una fracción mínima de las 144\,000
configuraciones. La búsqueda aleatoria y MMAS se ejecutaron en GPU
distintas, lo que se mitigó con semillas fijas, presupuesto en tokens y
la reevaluación de los mejores candidatos aleatorios en la misma GPU que
MMAS. El texto se normaliza a minúsculas, por lo que las métricas no
distinguen mayúsculas. El griego se entrenó sobre una submuestra de
300\,000 pares, y su filtrado conserva solo el 91\,\% de los pares, lo
que sugiere ruido de alineación.

% =====================================================================
\section{Conclusiones y trabajo futuro}

\pendiente{conclusiones}.

Como trabajo futuro se propone sustituir la discretización manual por
ACO$_\text{MV}$ \cite{liao2014acomv}, que trata de forma nativa variables
continuas, ordinales y categóricas; repetir la búsqueda con varias
semillas; ajustar directamente sobre el griego; y entrenar un modelo
bidireccional usando las etiquetas de dirección ya reservadas en los
tokenizadores.

% =====================================================================
\bibliographystyle{IEEEtran}
\bibliography{referencias}

\end{document}
